% !TEX root = ../ipeesec.tex
\section{Introduction}\label{sec:introduction}

Multiple ways exist to construct an encrypted tunnel between networks.
Common options are WireGuard~\cite{donenfeldWireGuard2017}, OpenVPN~\cite{openvpn} and IPsec~\cite{rfc4301}, where the last is the IETF standard and the oldest of them.
Through XFRM this protocol suite is implemented in the Linux kernel, and most routers, firewalls and operating systems support it.
In this paper, we wish to implement parts of this protocol using sheaf and type theory and formalize these findings as a Lean~4 implementation.

\Cref{fig:ipsec-overview} shows, on an abstract level, how two hosts set up and use an IPsec tunnel.
First, the \IKE{} protocol~\cite{rfc7296} runs in two exchanges.
In \texttt{IKE\_SA\_INIT} the initiator (client) and the responder (server) agree on cryptographic algorithms and perform a Diffie--Hellman key exchange. Both then derive the same keys, which form the \emph{IKE security association} (IKE~SA).
In \texttt{IKE\_AUTH}, already protected by these keys, both peers authenticate each other and negotiate a \emph{child SA}, consisting of the keys, algorithms and traffic selectors for the actual data traffic.
Each peer stores the child SA, one per direction and identified by its \SPI{}, in its security association database (\SAD{}).
Afterwards, every IP packet matched by the security policy database (\SPD{}) is encrypted and authenticated with \ESP{}~\cite{rfc4303} and sent through the tunnel.

\begin{figure}[htbp]
	\centering
	\begin{tikzpicture}[
			>=Stealth, font=\small,
			peer/.style={draw, thick, rounded corners=3pt, minimum width=2.6cm,
					minimum height=0.8cm},
			plain/.style={->, thin},
			ike/.style={->, thick},
			esp/.style={->, very thick, double, double distance=1.2pt},
			msg/.style={above, font=\footnotesize, inner sep=2pt},
			state/.style={draw, fill=white, rounded corners=2pt, font=\footnotesize,
					inner sep=3pt},
			phase/.style={decorate, decoration={brace, amplitude=5pt}, thick},
			phaselabel/.style={right, font=\footnotesize, xshift=6pt},
		]
		\def\R{8}      % x-position of the server
		\def\B{-6.6}   % y-position where the lifelines end

		% peers and lifelines
		\node[peer] (C) at (0, 0)  {Client};
		\node[peer] (S) at (\R, 0) {Server};
		\draw[densely dashed] (C.south) -- (0, \B);
		\draw[densely dashed] (S.south) -- (\R, \B);

		% 1. key exchange (unencrypted)
		\draw[plain] (0, -1.1) -- node[msg] {algorithms, key share} (\R, -1.1);
		\draw[plain] (\R, -1.8) -- node[msg] {algorithms, key share} (0, -1.8);
		\node[state] at (\R/2, -2.4) {both derive a shared key};

		% 2. authentication (encrypted with the shared key)
		\draw[ike] (0, -3.3) -- node[msg] {identity, proof} (\R, -3.3);
		\draw[ike] (\R, -4.0) -- node[msg] {identity, proof} (0, -4.0);

		% 3. data transfer through the tunnel
		\draw[esp] (0, -5.0) -- node[msg, yshift=2pt] {encrypted packet} (\R, -5.0);
		\draw[esp] (\R, -5.9) -- node[msg, yshift=2pt] {encrypted packet} (0, -5.9);

		% phases
		\draw[phase] (\R+1.6, -0.8) -- (\R+1.6, -2.7) node[midway, phaselabel] {key exchange};
		\draw[phase] (\R+1.6, -2.9) -- (\R+1.6, -4.3) node[midway, phaselabel] {authentication};
		\draw[phase] (\R+1.6, -4.5) -- (\R+1.6, -6.3) node[midway, phaselabel] {tunnel};
	\end{tikzpicture}
	\caption{Abstract overview of IPsec.
		Client and server first agree on a shared key, then authenticate each other, and finally exchange encrypted packets through the tunnel.
		Thin arrows are unencrypted, thick and double arrows are encrypted.}
	\label{fig:ipsec-overview}
\end{figure}
